\documentclass[a4paper, 12pt]{article}

\usepackage{cmap}
\usepackage[T2A]{fontenc}
\usepackage[english, russian]{babel}
\usepackage[utf8]{inputenc}
\usepackage[left=2cm,right=1.5cm,top=2cm,bottom=2cm]{geometry}
\usepackage{amsmath}
\usepackage{amssymb}
\usepackage{etoolbox}
\usepackage{amsthm}
\usepackage{booktabs}
\usepackage{graphicx}

\newcommand{\R}{\mathbb{R}}
\newcommand{\Q}{\mathbb{Q}}
\newcommand{\N}{\mathbb{N}}
\newcommand{\Z}{\mathbb{Z}}
\renewcommand{\phi}{\varphi}
\renewcommand{\epsilon}{\varepsilon}
\renewcommand{\emptyset}{\varnothing}
\newcommand{\lra}{\Leftrightarrow}
\renewcommand{\div}{\mid}
\newcommand{\ndiv}{\nmid}
\newcommand{\om}{\bar{\bar{o}}}
\newcommand\tab[1][.5cm]{\hspace*{#1}}

\theoremstyle{definition}
\newtheorem{definition}{Определение}[section]
\newtheorem{theorem}{Теорема}[section]
\newtheorem*{consequense}{Следствие}
\newtheorem{cons}{Следствие}
\newtheorem{lemma}{Лемма}[section]
\newtheorem{statement}{Утверждение}[section]
\newtheorem*{comm}{Замечание}
\newtheorem*{axiom}{Аксиома}
\newtheorem*{def20}{Определение 2.0}
\newtheorem*{algorithm}{Алгоритм}
\newtheorem*{introduction}{Введение}
\newtheorem*{eratosphene}{Решето Эратосфена}
\newtheorem*{example}{Пример}
\newtheorem*{examples}{Примеры}
\newtheorem*{theorem*}{Теорема}
\usepackage{titlesec}
\titleformat{\section}{\LARGE \bfseries}{\thesection}{1em}{}
\titleformat{\subsection}{\Large\bfseries}{\thesubsection}{1em}{}
\titleformat{\subsubsection}{\large\bfseries}{\thesubsubsection}{1em}{}

\usepackage{hyperref}
\usepackage{xcolor}
\definecolor{linkcolor}{HTML}{225ae2}
\definecolor{urlcolor}{HTML}{225ae2}
\hypersetup{
    pdfstartview=FitH, 
    linkcolor=linkcolor,
    urlcolor=urlcolor,
    colorlinks=true
}

\title{\textbf{Элементы теории чисел}}
\author{Лектор: Королев Максим Александрович}

\begin{document}
    \fontsize{14pt}{20pt}\selectfont
    \maketitle
    \vspace{1cm}
    \begin{center}
        \includegraphics[width=0.8\linewidth]{mehmat.png}
    \end{center}
    \vspace{1.5cm}
    \begin{center}
        Конспект: Кирилл Яковлев, 108 группа \tab[5.5cm] tg: @fourkenz
    \end{center}
    \newpage
    \tableofcontents
    \fontsize{14pt}{20pt}\selectfont
    \newpage
    \begin{introduction}
        Следующие понятия считаются интуитивно ясными:
    \end{introduction}
    \begin{enumerate}
        \item Понятие натурального ряда $\mathbb{N} = \{1, 2, 3, 4, 5, \dots\}$.
        \item У каждого натурального числа $n$ существует единственное натуральное число $m=n+1$ следующее за ним. 
        \item Понятие отрицательных чисел и нуля.
        \item Понятие суммы, разности и произведения двух целых чисел.
    \end{enumerate}
    \begin{axiom}
        Если $M \subset \mathbb{N}$ обладает следующими свойствами: $(1 \in M)$ и $(\forall n\in M$ выполнено $n+1 \in M)$, то $M = \mathbb{N}$.
    \end{axiom}
    \begin{cons}
        Всякое непустое подмножество натурального ряда содержит минимальный элемент.
    \end{cons}
    \begin{cons}
        Всякое непустое конечное подмножество натурального ряда содержит максимальный элемент.
    \end{cons}
    \begin{cons} (Принцип математической индукции)\newline
        Если известно, что некоторое утверждение о натуральных числах выполнено для натурального числа $a$, а также из предположения о том, что утверждение верно при некотором $n$ следует справедливость этого утверждения и для числа $n+1$, то это утверждение верно для всех натуральных чисел, больше или равных $a$.
    \end{cons}
    \section{Делимость целых чисел}
    \begin{definition}
        Пусть $a,b\in \mathbb{N}, b\ne 0$. Говорят что $a$ делится на $b$, если существует $c\in \mathbb{Z}$, такое, что $a=bc$.
    \end{definition}
    \begin{comm}
        $a$ называется делимым, а $b$ называется делителем числа $a$. Запись $b \mid a$ означает, что $b$ делит $a$. Если $b$ не делит $a$, то пишут $b\ndiv a$.
    \end{comm}
    \begin{lemma} \label{lemma1.1}
        Пусть $a,b,c \in \mathbb{Z}$, тогда:
        \begin{enumerate}
            \item $1 \mid a$.
            \item $a\ne 0 \Rightarrow a\mid a$.
            \item $a\mid b \Rightarrow a\mid bc$.
            \item $a\mid b$ и $b\mid c \Rightarrow a\mid c$.
            \item $a\mid b$ и $a\mid c \Rightarrow a\mid (b+c)$.
            \item $a\mid b$ и $b\ne 0 \Rightarrow |a|\leq|b|$.
        \end{enumerate}
    \end{lemma}
    \begin{theorem} \label{th1.1}
        Если $a\in \Z, b\in \N$, то существует единственная пара целых чисел $q$ и $r$, такие, что $a=bq+r$, где $0\leq r\leq b-1$. 
    \end{theorem}
    \begin{proof}
        Докажем существование: Если $a$ делится на $b$, то $a=bc$. В таком случае возьмем $q=c, r=0$. Теперь пусть $a$ не делится на $b$. Рассмотрим непустое множество $M$ натуральных чисел, представимых в виде $a-kb, k\in \mathbb{Z}$, возьмем $k=-(|a|+1)$, тогда $a-kb=b(|a|+1)+a\geq b(|a|+1)-|a|\geq 1\cdot (|a|+1)-|a|=1 \Rightarrow a-kb$ - натуральное. Значит, у $M$ есть минимальный элемент $a-kb$. Возьмем $q=k, r=a-kb=a-bq > 1$. Осталось показать, что $0\leq r\leq b-1$. Предположим, что $r\geq b$. Если $r=b$, то $a=bq+b=b(q+1)$ получаем противоречие, так как $a$ не делится на $b$. Значит, $r=b+m, m\geq 1$. Получаем $1\leq m=r-b<r$, при этом $a=bq+r=bq+b+m=b(q+1)+m \Rightarrow m=a-b(q+1) \Rightarrow m\in M$ и $m<r$, получаем противоречие, так как $a$ не делится на $b$. Доказано, что $r<b \Rightarrow$ представление $a=bq+r$ - искомое. Докажем единственность: предположим, что для некоторого $a$ и $b$ имеются пары чисел с указанным свойством: $q,r$ и $q_1,r_1$, причем $0\leq r\leq r_1 \leq b-1$. Тогда $a=bq+r=bq_1+r_1\Rightarrow 0\leq b(q-q_1)=r_1-r$. Значит, $b$ делит разность $r_1-r$. Допустим, что $q\ne q_1$, тогда по пункту 6 леммы \ref{lemma1.1} получаем $b\leq r_1-r$ и в то же время $r_1-r\leq b-1<b$. Получаем противоречие, поэтому $q=q_1$, следовательно, и $r=r_1$.
    \end{proof}
    \section{Наименьшее общее кратное и наибольший общий делитель (НОК и НОД)}
    \begin{def20}
         $n\geq 2, a_1, \dots, a_n \in \mathbb{N}$ пусть натуральное число $k$ делится на каждое из этих чисел. Тогда $k$ - общее кратное чисел $a_1, \dots, a_n$. \\
        Пусть $a_1, \dots, a_n$ - целые числа не все равные нулю. Натуральное число $d$ называется общим делителем $a_1, \dots, a_n$, если $d$ делит каждое из этих чисел.
    \end{def20}
    \begin{comm}
        Множество таких $k$ непусто, в нем лежит, например, произведение всех этих чисел. \\ Множество таких $d$ конечно: если $a_i\ne 0$, то $d$ находится среди делителей числа $a_i$, (по пункту 6 леммы \ref{lemma1.1}) $d\leq |a_i|$, значит числа $d$ образуют конечное множество, оно непусто, так как содержит единицу.
    \end{comm}
    \begin{definition}
        Наименьшее натуральное число, делящееся на каждое из чисел $a_1, \dots, a_n$, называют их наименьшим общим кратным, его обозначают $[a_1, \dots, a_n]$.
    \end{definition}
    \begin{theorem} \label{th2.1}
        Каждое общее кратное натуральных чисел $a_1, \dots, a_n$ делится на их НОК.
    \end{theorem}
    \begin{proof}
        Пусть $M$ - общее кратное $a_1, \dots, a_n$, $K=[a_1, \dots, a_n]$. Поделим $M$ на $K$ с остатком: $M=kq+r, 0\leq r\leq k-1\leq k$. Допустим, что $K\ne 0$. По определению, всякое число $a_i$ делит оба числа $M$ и $K \Rightarrow a_i$ делит разность $k=M-qK$, значит $k$ является общим кратным для $a_1, \dots, a_n$, но $k<K$, получаем противоречие, так как какое-то кратное оказалось меньше минимального. Значит, $k=0$ и $M=qK$.  
    \end{proof}
    \begin{definition}
        Наибольшее из натуральных чисел $d$, делящих каждое из чисел $a_1, \dots, a_n$, называют наибольшим общим делителем $a_1, \dots, a_n$, его обозначают $(a_1, \dots, a_n)$.
    \end{definition}
    \begin{definition}
        Числа $a$ и $b$ называется взаимно простыми, если $(a,b)=1$. Числа $a_1, \dots, a_n$ называются взаимно простыми в совокупности, если \\ $(a_1, \dots, a_n)=1$. Числа $a_1, \dots, a_n$ попарно взаимно просты, если $(a_i,a_j)=1$ $\forall i,j: 1\leq i<j\leq n$.
    \end{definition}
    \begin{theorem} \label{th2.2}
        $[a,b]\cdot (a,b)=ab, \forall a,b \in \mathbb{N}$.
    \end{theorem}
    \begin{proof}
        $ab$ - общее кратное $a$ и $b$. По теореме \ref{th2.1} $ab$ делится на $[a,b]$, то есть $ab=c[a,b]$, где $c\geq 1$ - натуральное число. Покажем, что $a$ и $b$ делятся на $c$. Действительно, $a=\frac{ab}{[a,b]}\cdot \frac{[a,b]}{b}=c\cdot \frac{[a,b]}{b}$, $b=\frac{ab}{[a,b]}\cdot \frac{[a,b]}{a}=c\cdot \frac{[a,b]}{a}$, но оба числа $\frac{[a,b]}{a}$ и $\frac{[a,b]}{b}$ - натуральные, значит $c$ - общий делитель $a$ и $b$. Пусть теперь $d$ - произвольный общий делитель $a$ и $b$, тогда $\frac{ab}{d}=a\cdot \frac{b}{d}$, то есть число $\frac{ab}{d}$ делится нацело на каждое из чисел $a$ и $b$. По теореме \ref{th2.1}, оно делится на $[a,b]$, то есть $\frac{ab}{d}=[a,b]m, $ где $ m\geq 1$ - натуральное число, но тогда $\frac{ab}{[a,b]}=c=dm$, то есть $d$ делит $c$. В силу пункта 6 леммы \ref{lemma1.1} $d\leq c$, значит $c=(a,b)$.
    \end{proof}
    \begin{theorem} \label{th2.3}
        Пусть $a,b,c \in \mathbb{N}$, причем $a \div bc$ и $(a,b)=1$, тогда $a \div c$. 
    \end{theorem}
    \begin{proof}
        $(a,b)=1 \Rightarrow$ (по теореме \ref{th2.2}) $bc$ делится нацело на $[a,b]=ab$, то есть $bc=abm$, где $m\geq 1$ - натуральное число. Сократим обе части на $b$, получим $c=am$. 
    \end{proof}
    \begin{theorem} \label{th2.4}
        Пусть $\Delta=(a,b)\geq 1 \Rightarrow (\frac{a}{\Delta}, \frac{b}{\Delta})=1$.
    \end{theorem} 
    \begin{proof}
        Пусть $m\in \N$ и $m \mid \frac{a}{\Delta}, m \mid \frac{b}{\Delta}$ предположим, что $m > 1 \Rightarrow cm=\frac{a}{\Delta},dm=\frac{b}{\Delta}\Rightarrow \Delta cm=a, \Delta dm=b \Rightarrow \Delta m \mid a$ и $\Delta m\mid b\Rightarrow \Delta m$ - общий делитель $a$ и $b$. Но т.к. $m>1$, то $\Delta m>\Delta \Rightarrow \Delta=(a,b) \leq \Delta m$ - противоречие, поскольку $\Delta$ - НОД $\Rightarrow m=1\Rightarrow (\frac{a}{\Delta}, \frac{b}{\Delta})=1$.
    \end{proof} 
    \section{Алгоритм Евклида}
    \begin{lemma} \label{lemma3.1}
        Пусть $a\in \Z, b\in \N$ и $b \mid a$. Тогда $(a,b)=b$.
    \end{lemma}
    \begin{proof}
        Пусть $(a,b)=c\Rightarrow c\mid b \Rightarrow$ (по лемме \ref{lemma1.1}) $c\leq b$, но $b\mid a,\\ b\mid b\Rightarrow b$ - общий делитель $a$ и $b\Rightarrow b\leq c\Rightarrow b=c=(a,b)$.
    \end{proof} 
    \begin{lemma} \label{lemma3.2}
        Пусть $a\in \Z, b\in \N, a=bq+r: r,q\in \Z, r\geq 0$. Тогда $(a,b)=(b,r)$. 
    \end{lemma} 
    \begin{proof}
        Пусть $\Delta=(a,b), \delta=(b,r)$. Имеем $\delta \div b \Rightarrow \delta \div bq, \delta \div r\Rightarrow$ (лемма \ref{lemma1.1}) $\delta \div bq+r=a\Rightarrow \delta \div a,\delta \div b\Rightarrow \delta$ - общий делитель $a$ и $b \Rightarrow \delta\leq \Delta$. \\ $\Delta \div b, \Delta \div bq, \Delta \div a\Rightarrow$ (лемма \ref{lemma1.1}) $\Delta \div a-bq=r\Rightarrow \Delta$ - общий делитель $b$ и $r\Rightarrow \Delta \leq \delta \Rightarrow \Delta = \delta$.
    \end{proof} 
    \begin{algorithm}
        Получаем, что при поиске НОД $a$ и $b, (a,b)$ можно заменять любой парой $(b,r)=(b,a-bq), q\in \Z$. Положим $r_0=a, r_1=b$. \\ 
        Выполняем деление с остатком:\\
        $$r_0=r_1q_1+r_2, 0<r_2<r_1\Rightarrow (r_0, r_1)=(r_1,r_2)$$
        $$r_1=r_2q_2+r_3, 0<r_3<r_2\Rightarrow (r_1, r_2)=(r_2,r_3)$$
        $$r_2=r_3q_3+r_4, 0<r_4<r_3\Rightarrow (r_2, r_3)=(r_3,r_4)$$
        \tab[8.5cm]\vdots 
        $$r_{n-2}=r_{n-1}q_{n-1}+r_n, 0<r_n<r_{n-1}\Rightarrow (r_{n-2}, r_{n-1})=(r_{n-1},r_n)$$
        $$r_{n-1}=r_nq_n$$
        $$\Rightarrow \text{(лемма \ref{lemma3.1})} (r_{n-1},r_n)=r_n\Rightarrow (a,b)=r_n$$
    \end{algorithm}
    \section{Решение в целых числах линейного уравнения с двумя неизвестными}
    Рассмотрим уравнение ($*$) $ax+by=c$, такое, что $a,b,c\in \Z, a$ и $b$ не равняются нулю одновременно. $x,y\in \Z$ - неизвестные.
    \begin{theorem} \label{th4.1}
        $(1)$ Уравнение $(*)$ разрешимо $\lra \Delta = (a,b) \div c$.\\
        $(2)$ В случае разрешимости, множество решений этого уравнения бесконечно, все решения имеют вид $x=x_0+\frac{b}{\Delta}t, y=y_0-\frac{a}{\Delta}t$, где $x_0,y_0$ - произвольное решение, а $t\in \Z$.
    \end{theorem} 
    \begin{proof}
        Докажем первый пункт:
        \begin{itemize}
            \item[$(\Rightarrow)$] Если $x,y$ - решение, то $\Delta \div ax, \Delta \div by\Rightarrow$ (лемма \ref{lemma1.1}) $\Delta \div ax+by\Rightarrow \Delta \div c$.
            \item[$(\Leftarrow)$] Не теряя общности, можем считать, что $a\geq b\geq 0$. Доказываем индукцией по сумме $a+b$. \\
            База: $a+b=1\Rightarrow b=0$ и $a=1 \Rightarrow$ уравнение имеет вид $ax=c\Rightarrow x=c$. Можем предъявить решение $x=c, y=0$. В этом случае $\Delta = (1,0) \div 1$. \\
            Шаг: $n\geq 1$, считаем, что утверждение доказано для всех уравнений с условием $a\geq b\geq 0$, $1\leq a+b\leq n$. Пусть $ax+by=c$, где $a\geq b\geq 0,\\ a+b=n+1$ и $\Delta = (a,b) \div c \Rightarrow$ докажем, что есть хотя бы одно решение. Пусть $b=0, ax=c, \Delta = (a,0)=a, a\div c\Rightarrow c=am\Rightarrow x=m, y=0$ - решение. Пусть $b\geq 1$. Рассмотрим уравнение $(a-b)X+bY=c,\\ a-b\geq 0, b\geq 1 >0$. $(a-b)+b=(a+b)-b=n+1-b\leq n$. $(a-b, b)=(a,b) \div c\\ \Rightarrow$ по предположению индукции есть целочисленное решение $X_0,X_0$. \\ $(a-b)X_0+bY_0=c\Rightarrow aX_0-b(Y_0-X_0) = c \Rightarrow x=X_0, y=Y_0-X_0$ - решение.\\
            Докажем второй пункт (проверим, что $x_0,y_0$ - решение):\\
            $a(x_0+\frac{b}{\Delta}t)+b(y_0-\frac{a}{\Delta}t)=ax_0+\frac{ab}{\Delta}t+ay_0-\frac{ab}{\Delta}t=ax_0+by_0$. Обратно: пусть $x_0,y_0$ и $x,y$ - различные решения. $ax_0+by_0=c, ax+by=c \\ \Rightarrow a(x-x_0)+b(y-y_0)=0\Rightarrow a(x-x_0)=b(y_0-y)$. $\Delta = (a,b) \\ \Rightarrow a=\alpha\Delta, b=\beta\Delta\Rightarrow$ (теорема \ref{th2.4}) $(\alpha,\beta)=1 \\ \Rightarrow \alpha\Delta(x-x_0)=\beta\Delta(y_0-y)\Rightarrow \alpha(x-x_0)=\beta(y_0-y) \\ \Rightarrow \alpha \div \beta(y_0-y)\Rightarrow \alpha \div (y_0-y)\Rightarrow y_0-y=\alpha t\Rightarrow \alpha (x-x_0)=\beta \alpha t\\ \Rightarrow x-x_0=\beta t$.
        \end{itemize}
    \end{proof}
    \section{Простые числа}
    \begin{definition}
        Натуральное число $n>1$ называется простым, если оно имеет ровно два делителя: $1$ и $n$. В противном случае это число называется составным.
    \end{definition}
    \begin{comm}
        Единица не причисляется ни к простым, ни к составным.
    \end{comm}
    \begin{lemma} \label{lemma5.1}
        Наименьший делитель натурального числа $n>1$, отличный от единицы - простое число.
    \end{lemma} 
    \begin{proof}
        Пусть $d\div n, 1 \textless d \leq n$, и $d$ - наименьший с этими свойствами. Пусть $d$ - составное. Тогда $\exists k: k\div d$ и $1<k<d$. По лемме \ref{lemma1.1} $k\div n$, но $1<k<d$ - противоречие с тем, что $d$ - минимальный.
    \end{proof} 
    \begin{theorem} \label{th5.1}
        Множество простых чисел бесконечно.
    \end{theorem}
    \begin{proof}
        Пусть множество простых конечно: $p_1<p_2<\dots<p_n$ - все простые числа. Рассмотрим число $N = p_1p_2\dots p_n+1$. По лемме \ref{lemma5.1} наименьший делитель $p>1$ числа $N$ - простое число. Но $p$ отлично от $p_1\dots p_k, p$ делит $N$ нацело, а $N$ при делении на каждое из $p_1\dots p_n$ дает остаток $1$ - противоречие.
    \end{proof} 
    Пусть $x>0$, через $\pi(x)$ обозначим количество простых чисел на отрезке $[0, x]$ ($\pi(x)$ - количество простых чисел не превосходящих $x$).
    $$\pi(x)=\sum\limits_{p\leq x}1$$
    (Теорема \ref{th5.1}) $\lra \pi(x)$ - не ограничена сверху $\lra \pi(x)\to +\infty$ при $x\to +\infty$.\\
    Гипотеза Лежандра: $\pi(x)=\frac{x}{\ln x-C}$, где $C=1,08366$. Позднее Гаусс выдвинет более сложное и более точное предположение.
    Из доказательства теоремы Чебышева: $\forall \epsilon >0\exists x_0=x_0(\epsilon)$, т.ч. $\forall x\geq x_0$ выполнено неравенство:
    $$(A-\epsilon)\frac{x}{\ln x}<\pi(x)<(B+\epsilon)\frac{x}{\ln x}$$
    $$A=\ln (\frac{2^{\frac{1}{2}} \ 3^{\frac{1}{3}} \ 5^{\frac{1}{5}}}{30^{\frac{1}{30}}}), B=\frac{6}{5}A$$
    Асимптотический закон распределения простых чисел:
    $$\lim\limits_{x\to +\infty}(\frac{\pi(x)}{\frac{x}{\ln x}})=1 \lra A=B=1 \lra \pi(x)=(1+\om(1))\frac{x}{\ln x}$$
    \begin{lemma} \label{lemma5.2}
        Пусть $N$ - составное число, $p$ - наименьший простой делитель. Тогда $p\leq \sqrt{N}$.
    \end{lemma} 
    \begin{proof}
        $N$ - составное $\Rightarrow \exists a,b: 1<a<N, 1<b<N$ и $ab=N$. Значит $a\div N, b\div N, p$ - наименьший $\Rightarrow p\leq a, p\leq b \Rightarrow p^2\leq ab=N\\ \Rightarrow p\leq \sqrt{N}$.
     \end{proof}
    \begin{eratosphene}
        Выписываем все числа от $2$ до $N$, первое число в таблице - простое, это 2. Вычеркиваем все числа кратные 2, кроме нее самой. Первое невычеркнутое число после 2 - это 3 - значит оно простое. Вычеркиваем все числа, кратные 3, кроме самой 3. Первое невычеркнутое число после 3 - простое и т.д. После того как найдено наибольшее простое $p$ не превосходящее $\sqrt{N}$ и вычеркнуты все числа кратные $p$, в таблице останутся лишь простые числа, не превосходящие $N$ и только они.
    \end{eratosphene}
    \begin{theorem} \label{th5.2}
        (Основная теорема арифметики)\\
        Каждое целое число, большее $1$, раскладывается в произведение простых чисел, притом единственным способом (с точностью до порядка сомножителей). 
    \end{theorem}
    \begin{proof}
        Существование:\\ Индукция по $n>1$. Числа $n=2, n=3$ - простые, для них это утверждение справедливо. Пусть $n>3$, и допустим, что справедливость утверждения проверена для всех $m<n$. Если $n$ - простое, то утверждение очевидно. Пусть $n$ - составное. По лемме \ref{lemma5.1} его наименьший делитель - простое число $\Rightarrow n=p_1k$, но $k=\frac{n}{p_1}\leq \frac{n}{2}<n$. По предположению индукции $k=p_2\dots p_r$, где $p_2,\dots,p_r$ - простые. $\Rightarrow n=p_1k=p_1p_2\dots p_r$ - искомое разложение.\\
        Единственность:\\
        Пусть $n=p_1\dots p_r=q_1\dots q_s,$ где $p_i,q_i$ - простые числа и $r\leq s$. Тогда\\ $p_1\dots p_r=q_1a_1$, где $a_1=q_2\dots q_s \Rightarrow p_1 \div q_1a_1$. Возможно два случая:\\
        1) $(p,q)>1 \Rightarrow p_1=q_1$.\\
        2) $(p,q)=1 \Rightarrow$ (теорема \ref{th2.3}) $p_1 \div a_1=q_2\dots q_s, a_1=q_2a_2, a_2=q_3\dots q_s,\\ p_1 \div q_2a_2\Rightarrow$ либо $p_1=q_2$, либо $p_1\div a_2$ и т.д. Но $a_1>a_2>\dots\geq 1 \Rightarrow$ на одном из шагов обязательно будет иметь место равенство $p_1=q_k, k\leq s$ (иначе оказалось бы, что $p_1 \div 1$, а это невозможно). Итак, $p_1$ совпадает с одним из чисел $q_1,\dots,q_s$. Будем считать, что $p_1=q_1 \Rightarrow p_2\dots p_r=q_2\dots q_s$ продолжаем рассуждение и получаем, что $p_2$ совпадает с одним из $q_2,\dots q_s$, пусть $p_2=q_2$ и т.д. Если $r<s$ после $r$ шагов получили бы противоречивое равенство: $1=q_{r+1}\dots q_s \\ \Rightarrow r=s$ и множества $\{p_1,\dots, p_r\}$ и $\{q_1, \dots , q_s\}$ совпадают.
    \end{proof}
    \begin{comm}
        $n>1, n=q_1\dots q_s \Rightarrow n$ можно записать в виде $n=p_1^{\alpha_1}\dots p_k^{\alpha_k}, \\ p_1<p_2<\dots<p_k$ - каноническое разложение $n$ на простые сомножители.
    \end{comm} 
    \begin{definition}
        Пусть $n=p_1^{\alpha_1}\dots p_k^{\alpha_k}, p$ - простое. Тогда
        $$\nu_p(n)=
        \begin{cases}
            0, \hspace{5pt} $если $p\ndiv n,\\
            \alpha, \hspace{5pt} $если$ p=p_i. 
        \end{cases}$$
    \end{definition}    
    \begin{lemma} (Свойства $\nu_p(n)$) \label{lemma5.3}
        \begin{enumerate}
            \item Для любых целых чисел $a,b>1$ и любого простого $p$ справедливо равенство: $\nu_p(ab)=\nu_p(a)+\nu_p(b)$.
            \item Пусть $m,n>1$ - целые числа, тогда $m \div n \lra \nu_p(m)\leq \nu_p(n)$ для любого простого $p$.
        \end{enumerate}
    \end{lemma}
    \begin{proof} \tab
        \begin{enumerate}
            \item При перемножении степеней с одинаковыми основаниями, их показатели складываются.
            \item    
            \begin{itemize}
                \item[$(\Rightarrow)$] Пусть $n=km \Rightarrow \nu_p(n)=\nu_p(k)+\nu_p(m)\geq \nu_p(m)$.
                \item[$(\Leftarrow)$] Все разности $\nu_p(n)-\nu_p(m)$ - целые неотрицательные. Рассмотрим число:
                $$k = \prod\limits_pp^{\nu_p(n)-\nu_p(m)}$$ 
                Если $k=1$, то $\nu_p(n)=\nu_p(m)$ для всех $p$ и $m=n$. В силу основной теоремы арифметики, в этом случае $m \div n$. Пусть $k>1$, тогда в силу пункта 1:
                $$km=\prod\limits_pp^{\nu_p(n)-\nu_p(m)}\cdot \prod\limits_pp^{\nu_p(m)}=\prod\limits_pp^{\nu_p(n)}=n$$
                то есть $m \div n$. 
            \end{itemize}
        \end{enumerate}
    \end{proof}
    \begin{lemma}
        Для любых $a,b\in \N$ справедливы равенства:
        $$[a,b]=\prod\limits_pp^{\max(\nu_p(a),\nu_p(b))}$$
        $$(a,b)=\prod\limits_pp^{\min(\nu_p(a),\nu_p(b))}$$
    \end{lemma}
    \begin{proof}
        Обозначим $K=[a,b],$ $N=\prod\limits_pp^{\max(\nu_p(a),\nu_p(b))}$ поскольку\\ $\nu_p(a)\leq \nu_p(N),$ $\nu_p(b)\leq \nu_p(N)$, то $a$ и $b$ делят $N$ в силу леммы \ref{lemma5.3}. Значит $N$ - общее кратное чисел $a$ и $b$. С другой стороны, поскольку $a$ и $b$ делят $K$, то по лемме \ref{lemma5.3} имеем $\nu_p(a)\leq \nu_p(K),\nu_p(b)\leq \nu_p(K)$, так что $\nu_p(K)\geq \max(\nu_p(a), \nu_p(b))=\nu_p(N)$ для любого простого $p$. Значит, $N\div K$, но $N\leq K \Rightarrow N=K$. Вторая часть утверждения следует из первой, если воспользоваться равенством
        $$(a,b)=\frac{ab}{[a,b]}$$
        и тем, что $x+y=\max(x,y)+\min(x,y)$ $\forall x,y\in \R$.
    \end{proof}
    \section{Мультипликативные функции}
    Обозначения и пояснения:
    \begin{enumerate}
        \item Обозначение $\sum\limits_{d\div n}f(d)$ - сумма значений функции $f$ по всем делиителям $d$ числа $n$.
        \item Двойная сумма вычисляется следующим образом: \[\sum\limits_{m=1}^M\sum\limits_{n=1}^Ng(mn)=\sum\limits_{n=1}^Ng(1,n)+\sum\limits_{n=1}^Ng(2,n)+\dots +\sum\limits_{n=1}^Ng(M,n)\]
    \end{enumerate}
    \begin{definition}
        Функция $f$, определенная на множестве $\N$ называется мультипликативной, если для любых взаимно простых $a,b\in \N$ выполнено равенство:
        $$f(ab)=f(a)f(b)$$
    \end{definition}
    \begin{theorem}(Простейшие свойтсва мультипликативных функций) \label{th6.1}\\
        Пусть $f,g$ - мультипликативные функции. Тогда:
        \begin{enumerate}
            \item Если $f\not\equiv 0$, то $f(1)=1$.
            \item Если $n=p_1^{\alpha_1}\dots p_r^{\alpha_r}$ - каноническое разложение $n$, то $f(n)=f(p_1^{\alpha_1})\dots f(p_r^{\alpha_r})$.
            \item Функция $h$, определенная для любого $n\in \N$ равенством\\ $h(n)=f(n)g(n)$ - мультипликативна. 
        \end{enumerate}
        \begin{proof}
            \begin{enumerate}
            \item Так как $f\not\equiv 0$, то $\exists a\in \N: f(a)\ne 0$. Тогда\\ $f(a)=f(a\cdot 1)=f(a)f(1) \Rightarrow f(1)=1$.
            \item 3. Напрямую следует из определения.
            \end{enumerate}
        \end{proof}
    \end{theorem}
    Для исседование дальнейших свойств мультипликативных функций потребуется несколько вспомогательных лемм
    \begin{lemma} \label{lemma6.1}
        Пусть $p$ - простое число, $r\geq 2$ и пусть целые числа $a_1,\dots, a_r$ попарно взаимно просты, причем $p \div a_1\dots a_r$. Тогда найдется номер $1\leq s\leq r$ такой, что $p\div a_s$.
    \end{lemma} 
    \begin{proof}
        Индукция по $r$. Если $r=2$, то это есть очевидно следствие теоремы \ref{th2.3}. Пусть $m\geq 3$ и утверждение доказано для всех $r\leq m-1$. Пусть $a_1,\dots a_m$ попарно взаимно просты и $p \div a_1\cdot...\cdot a_m$. Полагая $a=a_1\cdot ... \cdot a_{m-1}$ будем иметь: $p \div aa_m$. Если $(p,a)=1$, то $p\div a_m$ по теореме \ref{th2.3}. Пусть $(p,a)>1$. Так как $p$ - простое, то $(p,a)=p$ и $p$ делит некоторый сомножитель $a_{s}:\\ 1 \leq s \leq m-1$.
    \end{proof} 
    \begin{lemma} \label{lemma6.2}
        Пусть $b\div a$ и $c\div a$, причем $(b,c)=1$. Тогда $bc\div a$.
    \end{lemma}
    \begin{proof}
        Из условия следует, что $a$ - общее кратное $b$ и $c$. По теореме \ref{th2.1} $a$ делится на $[b,c]$, по теореме \ref{th2.2} $[b,c]=bc$.
    \end{proof}
    \begin{consequense}
        Пусть $r\geq 2$, и пусть целые числа $b_1\dots b_r$ попарно взаимно просты, причем $b_1 \div a, \dots, b_r \div a$. Тогда $b_1\dots b_r \div a$.
    \end{consequense} 
    \begin{proof}
        Индукция по $r$. Если $r=2$, получаем утверждение леммы. Пусть $m\geq 3$ и утверждение доказано для всех $r\leq m-1$. Пусть $b_1, \dots b_m$ попарно взаимно просты и каждое из них делит $a$. В силу предложения индукции, $a$ делится на произведение $b = b_1\dots b_{m-1}$. Заметим, что $(b,b_m)=1$. Действительно, в противном случае найдется простое число $p$, делящееся как на $b_m$ так и на $b$. По лемме \ref{lemma6.1} $p$ будет делить и некоторые $b_{\xi}: 1\geq \xi \geq m-1$. Следовательно $(b_m,b_{\xi})\geq p>1$, что противоречит условию. Так как $a$ делится на $b$ и $b_m$, и $(b, b_m)=1$, то в силу леммы \ref{lemma6.2} $a$ делится на $bb_m=b_1\dots b_m$.
    \end{proof}
    \begin{lemma} \label{lemma6.3}
        Пусть числа $a$ и $b$ взаимно просты, и пусть $d_1$ и $d_2$ пробегают соответственно множества всех делителей $a$ и $b$. Тогда величина $d=d_1d_2$ пробегает без повторений всё множество делителей числа $ab$.
    \end{lemma} 
    \begin{proof} \tab
        \begin{enumerate}
            \item Если $d_1 \div a,$ $d_2 \div b$, то $a=kd_1,$ $b=md_2$ при некоторых $k,m\in \Z$, так что $ab=kmd_1d_2$, то есть $d_1d_2$ - делитель $ab$.
            \item Допустим, что $d_1d_2=\delta_1\delta_2$ для некоторых чисел $d_1, \delta_1$ делящих $a$ и некоторых чисел $d_2, \delta_2$, делящих $b$. Очевидно, что $(d_1,\delta_2)=1$, так как в противном случае нашлось бы простое $p$, делящееся одновременно и $a$ и $b$, что невозможно. Но $d_1 \div \delta_1\delta_2$ по теореме \ref{th2.3} $d_1 \div \delta_1$ и, следовательно $d_1\leq \delta_1$. Аналогично доказывается, что $\delta_1\div d_1$ и $\delta_1\leq d_1$. Значит $d_1=\delta_1, d_2=\delta_2$, то есть все произведения $d_1$ и $d_2$ различны.
            \item Докажем, наконец, что всякий делитель $d$ числа $ab$ встретится среди произведений $d_1d_2$. Если $d=1$, то это очевидно. Пусть $d\geq 2$ и $p_1^{\alpha_1}\dots p_r^{\alpha_r}$ - каноническое разложение $d$. Число $q_1=p_1^{\alpha_1} \div ab$. Из теоремы \ref{th2.3} следует, что $q_1$ делит либо $a$, либо $b$ (но не оба сразу). То же верно и для чисел $q_{\xi}=p_{\xi}^{\alpha_{\xi}}, \xi=2,3,\dots,r$. Пусть, для определенности, $q_1,\dots, q_t$ - все сомножители, делящие $a$, и $q_{t+1},\dots, q_r$ - все сомножители, делящие $b$. По следствию леммы \ref{lemma6.2} произведение $d_1=q_1\dots q_t$ делит $a$, произведение $d_2=q_{t+1},\dots, q_r$ делит $b$, но $d_1d_2=d$.
        \end{enumerate}
    \end{proof}
    \begin{theorem} \label{th6.2}
        Пусть функция $f$ мультипликативна. Тогда функция $F$, определенная при любом $n\in \N$ равенством:
        $$F(n)=\sum\limits_{d\div n}f(d)$$
        мультипликативна.
    \end{theorem}
    \begin{proof}
        Пусть $(a,b)=1$. По лемме \ref{lemma6.3}, все делители $ab$ получим без повторений, рассмотрев все произведения $d=d_1d_2$, где $d_1 \div a$, $d_2 \div b$. Значит
        $$F(ab)=\sum\limits_{d\div ab}f(d)=\sum\limits_{d_1 \div a}\sum\limits_{d_2 \div b}f(d_1d_2)=\sum\limits_{d_1 \div a}\sum\limits_{d_2 \div b}f(d_1)f(d_2)=$$
        $$=(\sum\limits_{d_1 \div a}f(d_1))(\sum\limits_{d_2 \div b}f(d_2))=F(a)F(b).$$
        Взаимная простота $d_1$ и $d_2$ очевидна.
    \end{proof} 
    \begin{consequense}
        Если $p_1^{\alpha_1}\dots p_r^{\alpha_r}$ - каноническое разложение $n$, а $F$ - функция из условия теоремы, то
        $$F(n)=\prod\limits_{i=1}^{r}(1+f(p_i)+f(p_i^2)+\dots +f(p_i^{\alpha_i}))$$
        (при условии что $f\not\equiv 0$).
    \end{consequense}
    \begin{definition}
        Функция Мебиуса $\mu(n)$ определяется равенствами:
        \[\mu(n)=\begin{cases}
            1, \hspace{5pt} \text{если} \hspace{5pt} n=1,\\
            0, \hspace{5pt} \text{если} \hspace{5pt} n \hspace{5pt} \text{делится на квадрат простого числа},\\
            (-1)^k,\hspace{5pt} \text{если} \hspace{5pt} n=p_1\dots p_k - \text{различные простые числа}.
        \end{cases}\]
    \end{definition} 
    Примеры: $\mu(2)=(-1)^1=-1, \mu(3)=-1, \mu(4)=0, \mu(5)=-1, \mu(6)=(-1)^2=1, \mu(7)=-1, \mu(8)=\mu(9)=0, \mu(10)=(-1)^2=1$
    $n=p_1\dots p_k$, $m=q_1\dots q_r$, $(m,n)=1\Rightarrow mn=p_1\dots p_kq_1\dots q_r \Rightarrow \mu(mn)=(-1)^{k+r}=(-1)^k(-1)^r=\mu(m)\mu(n)$.
    $p$ - простое $\Rightarrow \mu(p)=-1, \mu(p^2)=0, \mu(p^3)=0, \dots$
    \begin{theorem}
        (Основное свойство функции Мебиуса) \label{th6.3}
        \[\sum\limits_{d\div n}\mu(d)=\begin{cases}
            1, \hspace{5pt} \text{если} \hspace{5pt} n=1,\\
            0, \hspace{5pt} \text{во всех остальных случаях}.
        \end{cases}\]
    \end{theorem} 
    \begin{proof}
        Пусть $F(n)=\sum\limits_{d\div n}\mu(d)\Rightarrow$ (По теореме \ref{th6.2}) $F$ - мультипликативна. Пусть $p$ - простое, $n=p^{\alpha}, \alpha\geq 1 \Rightarrow F(p^{\alpha})=\sum\limits_{d\div p^{\alpha}}\mu(d)=\mu(1)+\mu(p)+\mu(p^2)+\dots +\mu(p^{\alpha})=1-1=0$.
    \end{proof}
    \begin{definition}
        Функция Эйлера $\phi(n)$ определяется для натурального $n$ как количество чисел $m$ с условиями $1\leq m\leq n$, таких, что $(m,n)=1$
    \end{definition}
    \begin{examples}
        $\phi(1)=1$, $\phi(2)=1$, $\phi(3)=1+1+0=2$, $\phi(4)=1+0+1+0=2$, $\phi(5)=1+1+1+1+0=4$, $\phi(6)=1+0+0+0+1+0=2$
    \end{examples}
    \begin{theorem}
        Функция Эйлера $\phi$ мультипликативна. Кроме того, если $p_1,\dots,p_k$ - все различные делители $n$, тогда:
        \[\phi(n)=n\prod\limits_{p\div n}(1-\frac{1}{p})=n(1-\frac{1}{p_1})\dots(1-\frac{1}{p_k})\]
    \end{theorem} 
    \begin{proof}
        Надо подсчитать число тех $m$, для которых $(m,n)=1$. По теореме \ref{th6.3}
        \[\sum\limits_{d\div (m,n)}\mu(d)=\begin{cases}
            1, \hspace{5pt} \text{если} \hspace{5pt} (m,n)=1,\\
            0, \hspace{5pt} \text{иначе}.
        \end{cases}\]
        \[\Rightarrow \phi(n)=\sum\limits_{1\leq m\leq n}\ \sum\limits_{d\div (m,n)}\mu(d)=\sum_{d\div n}\mu(d)\sum\limits_{d\leq m\leq n, d\div m}1\]
        \[1\leq m=kd\leq n \Rightarrow 1\leq k\leq \frac{n}{d}\]
        \[\Rightarrow \sum_{d\div n}\mu(d)\sum\limits_{d\leq m\leq n, d\div m}1 =\sum\limits_{d\div n}\mu(d)\frac{n}{d}=n\sum\limits_{d\div n}\frac{\mu(d)}{d}\] Функции $\mu(d)$ и $\frac{1}{d}$ - мультипликативные $\Rightarrow \frac{\mu(d)}{d}$ - мультипликативна $\Rightarrow$ по теореме \ref{th6.2} $\Rightarrow \sum\limits_{d\div n}\frac{\mu(d)}{d}$ - мультипликативна $\Rightarrow \phi(n)$ - мультипликативна.\\
        $n=p^{\alpha},\ p$ - простое, $\alpha\geq 1$
        \begin{multline*}
            \Rightarrow \phi(p^{\alpha})=p^{\alpha}\sum\limits_{d \div p^{\alpha}}\frac{\mu(d)}{d}=p^{\alpha}(\frac{\mu(1)}{1}+\frac{\mu(p)}{p}+\frac{\mu(p^2)}{p^2}+\dots+\frac{\mu(p^{\alpha})}{p^{\alpha}})=\\
            =p^{\alpha}(1+\frac{\mu(p)}{p})=p^{\alpha}(1+\frac{1}{p})
        \end{multline*}.
        \begin{multline*}
            n=p_1^{\alpha_1}\dots p_k^{\alpha_k}\Rightarrow \phi(n)=\phi(p_1^{\alpha_1})\dots \phi(p_k^{\alpha_k})=p_1^{\alpha_1}(1-\frac{1}{p_1})\dots p_k^{\alpha_k}(1-\frac{1}{p_k})=\\
            =p_1^{\alpha_1}\dots p_k^{\alpha_k}(1-\frac{1}{p_1})\dots (1-\frac{1}{p_k})=n\prod\limits_{p\div n}(1-\frac{1}{p}).
        \end{multline*}
    \end{proof}
    \begin{theorem}(Формула обращения Мебиуса)\label{th6.5}
        Пусть $\forall n\geq 1$ функции $f$ и $g$ связаны соотношением
        \begin{equation}
            f(n)=\sum\limits_{d\div n}g(d)
        \end{equation}
        Тогда $\forall n\geq 1$ выполнено равенство
        \begin{equation}
            g(n)=\sum\limits_{d\div n}\mu(d)f(\frac{n}{d})
        \end{equation}
        Обратно, если $\forall n\geq 1$ $f$ и $g$ связаны соотношением $(2)$, то $\forall n\geq 1$ верно $(1)$. 
    \end{theorem} 
    \begin{proof}
        \begin{itemize}
            \item[$(\Rightarrow)$] Пусть выполнено $(1)$, преобразуем величину
            \begin{multline*}
                \sum\limits_{d\div n}\mu(d)f(\frac{n}{d})=\sum\limits_{d\div n}\mu(d)\sum\limits_{\delta\div \frac{n}{d}}g(\delta)=\sum\limits_{d\delta\div n}\mu(d)g(\delta)=\sum\limits_{\delta\div n}g(\delta)\sum\limits_{d\div \frac{n}{\delta}}\mu(d)=\\=(\text{по теореме } \ref{th6.3})\hspace{3pt} g(n).
            \end{multline*}
            Пояснение: 
            \[\sum\limits_{d\div \frac{n}{\delta}}\mu(d)=\begin{cases}
                1,\ \text{если}\ \frac{n}{\delta}=1,\\
                0,\ \text{если}\ \frac{n}{\delta}>1.
            \end{cases}\lra\ \begin{cases}
                1,\ \text{если}\ n=\delta,\\
                0,\ \text{если}\ n>\delta.
            \end{cases}\]
            \item[$(\Leftarrow)$] Пусть есть (2), преобразуем
            \[\sum\limits_{d\div n}\mu(d)f(\frac{n}{d})=\sum\limits_{d\div n}\sum\limits_{\delta\div d}\mu(\delta)f(\frac{d}{\delta})=\]
                \[(\delta\div d \Rightarrow d=\Delta\delta\Rightarrow \frac{d}{\delta}=\Delta)\]
            \[=\sum\limits_{\Delta\delta\div d}\mu(\delta)f(\Delta)=\sum\limits_{\Delta\div n}f(\Delta)\sum\limits_{\delta\div \frac{n}{\Delta}}\mu(\delta)=\text{(по теореме}\hspace{3pt} \ref{th6.3}) \hspace{3pt} f(n)\]
        \end{itemize}
    \end{proof} 
    \begin{consequense}
        \[\sum\limits_{d\div n}\phi(n)=n\]
    \end{consequense}
    \begin{proof}
        Выше доказали, что
        \[\phi(n)=\sum\limits_{d\div n}\mu(d)\frac{n}{d}\]
        это равенство (2), где $g(n)=\phi(n), f(k)=k$. По формуле обращения Мебиуса, для этих функций выполнено (1): $f(n)=n=\sum\limits_{d\div n}g(n)=\sum\limits_{d\div n}\phi(n)$ 
    \end{proof} 
    \begin{definition}
        Функция делителей $\tau(n)$ определяется, как число делителей натурального $n\geq 1$.
        \[\tau(n)=\sum\limits_{d\div n}1.\]
    \end{definition}
    \begin{comm}
        $f(1)\equiv 1$ - мультипликативна $\Rightarrow$ (по теореме \ref{th6.2}) $\tau(n)$ - мультипликативна.
    \end{comm} 
    \begin{statement}
        $n=p^{\alpha}, \ p$ - простое.
        \[\tau(p^{\alpha})=\sum\limits_{d\div p^{\alpha}}1=\alpha+1\]
        $n=p_1^{\alpha_1}\dots p_k^{\alpha_k}\Rightarrow \tau(n)=(\alpha_1+1)(\alpha_2+1)\dots(\alpha_k+1)$.
    \end{statement}
    \begin{definition}
        $\sigma(n)$ - сумма делителей числа $n\geq 1$
        \[\sigma(n)=\sum\limits_{d\div n}d\]
    \end{definition} 
    \begin{examples}
        $\sigma(6)=1+2+3+6=12$, $p$ - простое $\Rightarrow \sigma(p)=p+1$.
    \end{examples}
    Из теоремы \ref{th6.2} следует мультипликативность $\sigma(n)$.\\
    $n=p^{\alpha} \Rightarrow \sigma = 1+p+p^2+\dots p^{\alpha}=\frac{p^{\alpha+1}-1}{p-1}$.\\
    Если $n=p_1^{\alpha_1}\dots p_k^{\alpha_k} \Rightarrow \sigma(n)=\sigma(p_1^{\alpha_1})\dots \sigma(p_k^{\alpha_k})=\prod\limits_{s=1}^r\frac{p_1^{\alpha_s+1}-1}{p_s-1}=\prod\limits_{p^{\alpha} \div\div n}\frac{p^{\alpha+1}-1}{p-1}$
    \begin{comm}
        Функции $\tau(n)$ и $\sigma(n)$ - частный случай функции $\sigma_{\beta}(n)$, $\beta$ - любое вещественное число.
        $\sigma_{\beta}(n)\sum\limits_{d\div n}1=\tau(n)$, $\sigma_1(n)=\sigma(n)$.
    \end{comm}
    Упражнение: Доказать, что $\sigma(n)+\phi(n)=n\tau(n)$ имеет место $\lra n$ - простое.
    \begin{definition}
        Делитель $d$ числа $n$ называется собстввенным, если $d<n$.
    \end{definition} 
    \begin{definition}
        Число $n$ называется совершенным, если оно равно сумме своих собственных делителей: $n=\sigma(n)-n\lra \sigma(n)=2n$
    \end{definition}
    \begin{examples}
        $\sigma(6)=12=6\cdot 2,\ \sigma(28)=56=2\cdot 28$.
    \end{examples}
    \begin{theorem} (Эйлер)\\
        Четное число является совершенным $\lra$ когда оно имеет вид $2^{p-1}(2^p-1)$, где $p$ и $2^p-1$ - простые числа. (без доказательства)
    \end{theorem} 
    Простые числа вида $M_p=2^p-1$, где $p$ - простое, называются простыми Мерсена. Сейчас известно 51 простое число Мерсена. Самое большое из них отвечает простому $p=82589933$. В записи $M_p$ - 24862048 цифр. (результат получен 21.12.2018) Неизвестно, конечно или нет множество простых Мерсена. Гипотеза: если $\pi_M(x)$ - число простых Мерсена не превосходящих $x$, то\\
    $\pi_M(x)\thickapprox \ln\ln x$. Неизвестно, существуют или нет нечетные совершенные числа. Если $N$ - нечетное совершенное число, то
    \begin{itemize}
        \item[(1)] $N>10^{1500}$ (2012г.)
        \item[(2)] Наибольший простой делитель $N$ превосходит $10^8$ (2008г.)
        \item[(3)] Второй по величине простой делитель $N$ превосходит $10^4$ (1999г.)
        \item[(4)] Пусть $k\geq 1$. Тогда имеется не более чем $2^{4^k}$ несчетных совершенных чисел, имеющих ровно $k$ различных простых делителей. (2003г.)  
    \end{itemize}
    \begin{definition}
        Числа $a$ и $b\ (1<a<b)$ называются дружественными, если $(a)$ $a$ есть сумма собственных делителей $b$, $(b)$ число $b$ - сумма собственных делителей $a$:
         $\ \begin{cases}
            \sigma(b)-b=a,\\
            \sigma(a)-a=b.
         \end{cases} \lra\ \ \ \sigma(a)=\sigma(b)=a+b$.   
    \end{definition}
    \begin{examples}
        (ЕЩЕ НЕ ГОТОВО)
    \end{examples}
    Неизвестно, конечно или нет множество дружественных пар чисел. Сейчас извество $1229319267$ таких пар.
    Пусть $A(x)$ - число дружественных пар с $a\leq x$. $\frac{A(X)}{x}\to 0$ при $x\to \infty$ (П. Эрдеш 1955г.)
    \newpage
    \section{Непрерывные дроби}
    \begin{example}
        Заметим, что  $43=19\cdot 2+5$, $19=5\cdot 3+4$. Рассмотрим дробь:
        \begin{multline*}
            \cfrac{a}{b}=\cfrac{19}{43}=\cfrac{1}{\cfrac{43}{19}}=\cfrac{1}{\cfrac{2\cdot 19+5}{19}}=\cfrac{1}{2+\cfrac{5}{19}}=\cfrac{1}{2+\cfrac{1}{\cfrac{19}{5}}}=\cfrac{1}{2+\cfrac{1}{\cfrac{5\cdot 3+4}{5}}}=\\
            =\cfrac{1}{2+\cfrac{1}{3+\cfrac{1}{\cfrac{5}{4}}}}=\cfrac{1}{2+\cfrac{1}{3+\cfrac{1}{1+\cfrac{1}{4}}}}
        \end{multline*}
    \end{example}
    \begin{definition}
        Непрерывной (цепной) дробью будем называть выражение вида:
        \[[q_0; q_1, q_2, \dots q_n] = q_0=\cfrac{1}{q_1+\cfrac{1}{q_2+\cfrac{1}{\ddots+\cfrac{1}{q_n}}}}\ \ \ (\ast)\]
    \end{definition}
    \begin{theorem} \label{th7.1}
        Пусть $a$ - целое, $b$ - натуральное и пусть $(a,b)=1$. Пусть кроме того, $q_0, q_1,\dots, q_n$ - все неполные частные, возникающие при отыскании $(a,b)$ с помощью алгоритма Евклида. Тогда число $\alpha=\frac{a}{b}$ разлагается в непрерывную дробь вида ($\ast$).
    \end{theorem} 
    \begin{proof}
        Доказательство следует из цепочки равненств:\\
        \[a=bq_0+r_1,\]
        \[b=r_1q_1+r_2,\]
        \[r_1=r_2q_2+r_3,\]
        \[\vdots\]
        \[r_{n-2}=r_{n-1}q_{n-1}+r_n\]
        Получаем:
        \begin{equation*}
            \cfrac{a}{b}=q_0+\cfrac{r_1}{b}=q_0+\cfrac{1}{\cfrac{b}{r_1}}=q_0+\cfrac{1}{q_1+\cfrac{r_2}{r_1}}=q_0+\cfrac{1}{q_1+\cfrac{1}{\cfrac{r_1}{r_2}}}=q_0+\cfrac{1}{q_1+\cfrac{1}{q_2+\cfrac{r_3}{r_2}}}=\dots
        \end{equation*}
        Собирая полученые равенства вместе приходим к $(\ast)$.
    \end{proof}
    \begin{example}\tab \\
        $a=37,\ b=8,\\
        37=8\cdot 4+5,\\
        8=5\cdot 1+3,\\
        5=3\cdot 1+2, \tab[1cm] \Rightarrow \alpha=\frac{37}{8}=[4;1,1,1,2].\\
        3=2\cdot 1+1,\\
        2=1\cdot 2$.
    \end{example}
    \begin{definition}
        Величины $q_0, q_1, \dots, q_n$ в разложении числа $\alpha=\frac{a}{b}$ из теоремы \ref{th7.1} называется неполным частным $b$ в разложении $\alpha$ в непрерывную дробь.
    \end{definition} 
        Дроби \[\delta_0=q_0\]
        \[\delta_1=q_0+\cfrac{1}{q_1}\] 
        \[\delta_2=q_0+\cfrac{1}{q_1+\cfrac{1}{q_2}}\]
        \[\vdots\]
    называются подходящими дробями.
    \begin{example}
    \[q_0=4,\ q_1=1,\ q_2=1,\ q_3=1, q_4=2.\]
    Тогда
    \[\delta_0=4\]
    \[\delta_1=4+\cfrac{1}{1}=5\]
    \[\delta_2=4+\cfrac{1}{1+\cfrac{1}{1}}=4+\cfrac{1}{2}=\cfrac{9}{2}\]
    \[\delta_3=4+\cfrac{1}{1+\cfrac{1}{1+\cfrac{1}{1}}}=4+\cfrac{2}{3}=\cfrac{14}{3}\]
    \end{example}
    Пусть $\alpha$ - число, не являющееся рациональным (такие числа будем называть иррациональными). Тогда для $\alpha$ тоже можно построить разложение в непрерывную дробь. Это разложение будет бесконечным (в отличии от рационального $\alpha=\frac{a}{b}$), поэтому такое построение требует определенной аккуратности и проводится в несколько шагов. На первом шаге строятся подходящие дроби, отвечающие числу $\alpha$, затем исследуются их свойства. В итоге доказывается сходимость последовательности подходящих дробей к числу $\alpha$, что и завершает построение.\\
    Этап первый:\\
    Определим целое $q_0$ так, чтобы выполнялись неравенства: \[q_0<\alpha<q_0+1\] и положим $\alpha_0=\alpha$, так что \[q_0<\alpha_0<q_0+1\] но тогда $\alpha_0=q_0+\beta_0$, где $0<\beta_0<1$ и, следовательно, \[\alpha_1=\frac{1}{\beta_0}>1\] и \[\alpha_0=q_0+\frac{1}{\alpha_1}\] Число $\alpha_1$, очевидно, иррационально, определим по нему целое число $q_1$ так, чтобы выполнялись неравенства: \[q_1<\alpha_1<q_1+1\] но $\alpha_1>1$, так что $q_1\geq 1$, т.е. $q_1$ - натуральное. Далее \[\alpha_1=q_1+\beta_1\] где $0<\beta_1<1$ и следовательно \[\alpha_2=\frac{1}{\beta_1}>1,\ \alpha_1=q_1+\frac{1}{\alpha_2},\ \alpha_0=q_0+\cfrac{1}{q_1+\cfrac{1}{\alpha_2}}\] Число $\alpha_2$ также иррационально. Повторяя это процесс далее, получим бесконечные последовательности иррациональных чисел $\alpha_1, \alpha_2,\dots, \alpha_{\xi},\dots$ (причем $\alpha_{\xi}>1$ для всех $\xi$) и натуральных чисел $q_1,\dots, q_{\xi}$ таких, что \[q_{\xi}=[\alpha_{\xi}]\ \text{и} \ \alpha_{\xi}=q_{\xi}+\frac{1}{\alpha_{\xi}+1}\] 
    Величины $q_0,q_1,q_2,\dots$ станем называть неполными частными разложения $\alpha$ в непрерывную дробь.
    Несложно видеть, что при любом $\xi$ справедливо равенство \[\alpha=q_0+\cfrac{1}{q_1+\cfrac{1}{q_2+\cfrac{1}{\ddots+\cfrac{1}{q_{\xi-1}+\cfrac{1}{\alpha_{\xi}}}}}}\]
    Определим по этим числам последовательность подходящих дробей $\delta_{\xi},\\
    \xi=0,1,\dots$ равенствами
    \[\delta_0=q_0,\ \ \delta_1=q_0+\cfrac{1}{q_1},\ \ \delta_2=q_0+\cfrac{1}{q_1+\cfrac{1}{q_2}},\ \ \dots\]
    \[\delta_{\xi}=q_0+\cfrac{1}{q_1+\cfrac{1}{q_2+\cfrac{1}{\ddots+\cfrac{1}{q_{\xi-1}+\cfrac{1}{\alpha_{\xi}}}}}}\]
    Выпишем первые три такие дроби:
    \[\delta_0=q_0,\ \ \delta_1=\cfrac{q_0 q_1+1}{q_1},\ \ \delta_2=\cfrac{q_0 q_1 q_2+q_0+q_2}{q_1 q_2+1}\]
    обозначим их еще так:
    \[\delta_0=\frac{P_0}{Q_0},\ \text{где}\ \ P_0=q_0,\ \ Q_0=1\]
    \[\delta_1=\frac{P_1}{Q_1},\ \text{где}\ \ P_1=q_0 q_1+1,\ \ Q_1=q_1\]
    \[\delta_2=\frac{P_2}{Q_2},\ \text{где}\ \ P_2=q_0 q_1 q_2+1,\ \ Q_2=q_1 q_2+1\]
    посмотрим как эти величины связаны между собой:
    \[P_1=q_1 P_0+1,\ \ Q_1=q_1 Q_0+0\]
    \[P_2=q_2 (q_0q_1+1)+q_0=q_2 P_1+P_0,\ \ Q_2=q_2 Q_1+Q_0\]
    Введем (формально) величины $P_{(-1)}=1, Q_{(-1)}=0$ (к подходящим дробям они не имеют отношения: выражение $\frac{P_{(-1)}}{Q_{(-1)}}=\frac{1}{0}$ не определено)\\
    Тогда равенства для $P_1, Q_1, P_2, Q_2$ запишутся единообразно:
    \[P_{\xi}=q_{\xi}P_{\xi-1}+P_{\xi-2},\ \ Q_{\xi}=q_{\xi}Q_{\xi-1}+Q_{\xi-2}\ \ \ (\xi=1,2)\]
    Оказывается, эти соотношения верны и для всех $\xi\geq 3$. Чтобы аккуратно доказать их поступим следующим образом.
    Этап второй:\\
    Пусть даны переменные $x_0, x_1, x_2,\dots x_{\xi}, \dots$ произвольной природы (не обязательно целые числа). Рассмотрим величины $P_{\xi}$ и $Q_{\xi}$, определенные рекуррентными соотношениями
    \begin{equation} \label{equation3}
        \begin{cases}
            P_{\xi}=x_{\xi} P_{\xi-1}+P_{\xi-2}\\
            Q_{\xi}=x_{\xi} Q_{\xi-1}+Q_{\xi-2}
        \end{cases}
    \end{equation}
    ясно, что $P_{\xi}$ и $Q_{\xi}$ - некоторые многочлены от переменных $x_1,\dots x_{\xi},\dots$, например: $P_3-x_3P_2+P_1=x_3(x_0x_1x_2+x_0+x_2)+x_0x_1+1=x_0x_1x_2x_3+x_0x_1+x_0x_3+x_2x_3+1$, положим также $h_{\xi}=P_{\xi}Q_{\xi-1}-P_{\xi-1}Q_{\xi}$.
    \begin{lemma}\label{lemma7.1}
        При любом $\xi\geq 0$ справедливо равенство: $h_{\xi}=(-1)^{\xi-1}$
    \end{lemma} 
    \begin{proof}
        Индукция по $\xi$. В случае $\xi=0$ имеем: \[h_0=P_0Q_{(-1)}-P_{(-1)}Q_0=-P_{(-1)}Q_0=-1=(-1)^{0-1}\]
        Пусть соотношение доказано для всех $\xi\leq m$. Тогда 
        \begin{multline*}
            h_{m+1}=P_{m+1}Q_m-P_m Q_{m+1}=(x_{m+1}P_m+P_{m-1})Q_m-P_m(x_{m+1}Q_m+Q_{m-1})=\\=x_{m+1}(P_m Q_m-P_m Q_m)+P_{m-1}Q_m-P_m Q_{m-1}=-h_m=-(-1)^{m-1}=(-1)^m
        \end{multline*}
    \end{proof} 
    \begin{lemma}\label{lemma7.2}
        Если $x_0=q_0,\ x_1=q_1,\ x_{\xi}=q_{\xi}$ - целые числа, а величины $P_{\xi}$ и $Q_{\xi}$ определены в (3) то справедливы равентсва \[(P_{\xi},Q_{\xi})=(P_{\xi},P_{\xi-1})=(Q_{\xi},Q_{\xi-1})=1\]
    \end{lemma} 
    \begin{proof}
        Сразу следует из леммы \ref{lemma7.1}.
    \end{proof} 
    \begin{lemma}\label{lemma7.3}
        Пусть $x_1,\dots x_{\xi},\dots$ - произвольные переменные, и пусть выражения $\Delta_0,\dots \Delta_{\xi}\dots$ зависящие от $x_1,\dots x_{\xi},\dots$ определяются следующим образом: $\Delta_0=x_0$, а при $\xi\geq 1$ выражение для $\Delta_{\xi}$ получим заменив в выражении для $\Delta_{\xi-1}\ \ x_{\xi-1}$ на $x_{\xi-1}+\frac{1}{x_{\xi}}$, так что, например,
        \[\Delta_1=x_0+\cfrac{1}{x_1},\ \ \Delta_2=x_0+\frac{1}{x_1+\cfrac{1}{x_2}},\ \ \Delta_3=x_0+\frac{1}{x_1+\cfrac{1}{x_2+\cfrac{1}{x_3}}}\]
        тогда при любом $\xi\geq 0$ справедливо равенство $\Delta_{\xi}=\frac{P_{\xi}}{Q_{\xi}}$, где P и Q определены соотношениями \eqref{equation3}
    \end{lemma} 
    \begin{proof}
        Индукция по $\xi$. В случае $\xi=0,1$ эти соотношения фактическибыли проверены ранее. Пусть они верны для всех $\xi\leq m$. Тогда
        \[\Delta_{\xi}=\frac{P_m}{Q_m}=\frac{x_mP_{m-1+P_{m-2}}}{x_mQ_{m-1}+Q{m-2}}\]
        по определению, $\Delta_{m+1}$ получим из $\Delta_m$ заменой $x_m$ на $x_m+\frac{1}{x_{m+1}}$ переменная $x_m$, очевидно не входит в выражения для $P_{m-1}, p_{m-2}, Q_{m-1}, Q_{m-2}$. Следовательно
        \begin{multline*}
        \Delta_{m+1}=\cfrac{(x_m+\cfrac{1}{x_{m+1}})P_{m-1}+P_{m-2}}{(x_m+\cfrac{1}{x_{m+1}})Q_{m-1}+Q_{m-2}}=\cfrac{(x_{m+1}x_m+1)P_{m-1}+x_{m-1}P_{m-2}}{(x_{m+1}x_m+1)Q_{m-1}+x_{m-1}Q_{m-2}}=\\=\cfrac{x_{m+1}(x_m P_{m-1}+P_{m-2})+P_{m-1}}{x_{m+1}(x_m Q_{m-1}+Q_{m-2})+Q_{m-1}}=\cfrac{x_{m+1}P_m+P_{m-1}}{x_{m+1}Q_m+Q_{m-1}}
        \end{multline*}
        но числитель и знаменатель последней дроби совпадают в силу \eqref{equation3} с $P_{m+1}$ и $Q_{m+1}$
    \end{proof} 
    \begin{theorem} \label{th7.2}
        Пусть $\alpha$ - произвольное вещественное число, и пусть $q_0, q_1, \dots$ - конечная или бесконечная последовательность неполных частных разложения $\alpha$ в непрерывную дробь. Тогда подходящие дроби $\delta_{\xi}, \ \xi=0,1,\dots$, отвечающие такому разложению, вычисляются по формулам
        \begin{equation} \label{equation4}
            \delta_{\xi}=\frac{P_{\xi}}{Q_{\xi}}
        \end{equation}
        где величины $P_{\xi}$ и $Q_{\xi}$ определяются следующими реккурентными соотношениями: $P_{\xi}=q_{\xi}P_{\xi-1}+P_{\xi-2}, \ Q_{\xi}=q_{\xi}Q_{\xi-1}+Q_{\xi-2}$ с начальными условиями $P_{(-1)}=1,\ Q_{(-1)}=0,\ P_0=q_0,\ Q_0=1$. Все дроби \eqref{equation4} при этом несократимы.
    \end{theorem} 
    \begin{proof}
        Равенство \eqref{equation4} есть прямое следствие леммы \ref{lemma7.1}
    \end{proof} 
    Если $\alpha \not\in \Q \Rightarrow q_0, q_1, q_2, \dots\ , \delta_{\xi}=\frac{P_{\xi}}{Q_{\xi}}$. Осталось непонятным какое отношение имеют дроби $\delta_{\xi}$ к числу $\alpha$.\\
    Этап третий:\\
    \begin{lemma} \label{lemma7.4}
        При любом $\xi\geq 1$ верны неравенства: $\delta_{2\xi}>\delta_{2\xi-2}\ (\delta_{2\xi+1}<\delta_{2\xi-1})$ то есть подходящие дроби с четными (нечетными) номерами образуют монотонно возрастающую (убывающую последовательность).
    \end{lemma} 
    \begin{proof}
        \begin{multline*}
        \delta_{k}-\delta_{k-1}=\frac{P_k}{Q_k}-\frac{p_{k-1}}{Q_{k-1}}=\frac{P_k Q_{k-1}-P_{k-1}Q_k}{Q_k Q_{k-1}}= \text{(по лемме \ref{lemma7.1})}\\ =\frac{h_k}{Q_k Q_{k-1}}=\frac{(-1)^{k-1}}{Q_k Q_{k-1}}
        \end{multline*}
        тогда
        \begin{multline*}
            \delta_{2\xi}-\delta_{2\xi-2}=(\delta_{2\xi}-\delta_{2\xi-1})+(\delta_{2\xi-1}-\delta_{2\xi-2})=\frac{(-1)^{2\xi-1}}{Q_{2\xi}Q_{2\xi-1}} + \frac{(-1)^{2\xi-2}}{Q_{2\xi-1}Q_{2\xi-2}}=\\
            =\frac{1}{Q_{2\xi-1}}(\frac{1}{Q_{2\xi-2}}-\frac{1}{Q_{2\xi}})=\frac{Q_{2\xi}-Q_{2\xi-2}}{Q_{2\xi} Q_{2\xi-1} Q_{2\xi-2}}= (Q_{2\xi}=q_{2\xi} Q_{2\xi-1}+Q_{2\xi-2})\\
            =\frac{q_{2\xi} Q_{2\xi-1}}{Q_{2\xi} Q_{2\xi-1} Q_{2\xi-2}}=\frac{q_{2\xi}}{Q_{2\xi}Q_{2\xi-2}}>0
        \end{multline*}
        Неравенство $\delta_{2\xi+1}-\delta{2\xi-1}$ доказывается аналогично.
    \end{proof} 
    \begin{lemma} \label{lemma7.5}
        В условиях леммы \ref{lemma7.4} справедливы неравенства: $\delta_{\xi}<\alpha,\ \xi$ - четное и $\delta_{\xi}>\alpha,\ \xi$ - нечетное.
    \end{lemma} 
    \begin{proof}
        Рассмотрим выражения
        \[\alpha=q_0+\cfrac{1}{q_1+\cfrac{1}{q_2+\cfrac{1}{\ddots+\cfrac{1}{q_{\xi}+\cfrac{1}{\alpha_{\xi+1}}}}}}\]
        \[\delta_{\xi+1}=q_0+\cfrac{1}{q_1+\cfrac{1}{q_2+\cfrac{1}{\ddots+\cfrac{1}{q_{\xi}+\cfrac{1}{q_{\xi+1}}}}}},\ \ \delta_{\xi}=q_0+\cfrac{1}{q_1+\cfrac{1}{q_2+\cfrac{1}{\ddots+\cfrac{1}{q_{\xi}}}}}\]
        выражения для $\alpha$ и $\delta_{\xi+1}$ получаются из выражения для $\delta_{\xi}$ формальной заменой $q_{\xi}$ на $q_{\xi}+\frac{1}{\alpha_{\xi+1}}$ и на $q_{\xi}+\frac{1}{q_{\xi+1}}$ соответственно.
        \[\delta_{\xi}=\frac{P_{\xi}}{Q_{\xi}}=\frac{q_{\xi}P_{\xi-1}+P_{\xi-2}}{q_{\xi}Q_{\xi-1}+Q_{\xi-2}} \Rightarrow \alpha=\cfrac{(q_{\xi}+\cfrac{1}{\alpha_{\xi+1}})P_{\xi-1}+P_{\xi-2}}{(q_{\xi}+\cfrac{1}{\alpha_{\xi+1}})Q_{\xi-1}+Q_{\xi-2}}=\frac{A_{\xi}}{B_{\xi}}\]
        \[\delta_{\xi+1}=\cfrac{(q_{\xi}+\cfrac{1}{q_{\xi+1}})P_{\xi-1}+P_{\xi-2}}{(q_{\xi}+\cfrac{1}{q_{\xi+1}})Q_{\xi-1}+Q_{\xi-2}}=\frac{P_{\xi+1}}{Q_{\xi+1}}\]
        вычислим:
        \[\alpha-\delta_{\xi+1}=\frac{A_{\xi}}{B_{\xi}}-\frac{P_{\xi+1}}{Q_{\xi+1}}=\frac{A_{\xi}Q_{\xi+1}-B_{\xi}P_{\xi+1}}{B_{\xi}Q_{\xi+1}}\]
        числитель:
        \begin{multline*}
            \ast = (q_{\xi}+\frac{1}{q_{\xi+1}})(q_{\xi}+\frac{1}{\alpha_{\xi+1}})(P_{\xi-1}Q_{\xi-1}-P_{\xi-1}Q_{\xi-1})+(P_{\xi-2}Q_{\xi-2}-P_{\xi-2}Q_{\xi-2})+\\
            +(q_{\xi}+\frac{1}{q_{\xi+1}})+(P_{\xi-2}Q_{\xi-1}-P_{\xi-1}Q_{\xi-2})+(q_{\xi}+\frac{1}{\alpha_{\xi+1}})+(P_{\xi-1}Q_{\xi-2}-P_{\xi-2}Q_{\xi-1})
        \end{multline*}
        итак, числитель разности $\alpha-\delta_{\xi+1}$ равен
        \begin{multline*}
            (P_{\xi-1}Q_{\xi-2}-P_{\xi-2}Q_{\xi-1})(q_{\xi}+\frac{1}{\alpha_{\xi+1}}-q_{\xi}-\frac{1}{q_{\xi+1}})=h_{\xi-1}(\frac{1}{\alpha_{\xi+1}}-\frac{1}{q_{\xi+1}})=\\
            =(-1)^{\xi-2}\frac{q_{\xi+1}-\alpha_{\xi+1}}{\alpha_{\xi+1}q_{\xi+1}}=\frac{(-1)^{\xi+1}\alpha_{\xi+1}}{\alpha_{\xi+1}q_{\xi+1}}
        \end{multline*}
        $\Rightarrow$ знак разности $\alpha-\delta_r$ совпадает с $(-1)^r$
    \end{proof} 
    \begin{theorem}\label{th7.3}
        Последовательность подходящих дробей, отвечающих разложению иррационального числа $\alpha$ в непрерывную дробь, сходится к числу $\alpha$.
    \end{theorem} 
    \begin{proof}
        По лемме \ref{lemma7.4} последовательность $\delta_{2\xi},\ \xi=0,1,2,\dots$ монотонно возрастает. По лемме \ref{lemma7.5} она ограничена сверху числом $\alpha$. Аналогично, последовательность $\delta_{2\xi+1}$ монотонно убывает и ограничена снизу числом $\alpha$. По известной теореме из математического анализа, эти последовательности имеют пределы. По лемме \ref{lemma7.5}\ : $\delta_{2\xi}<\alpha<\delta_{2\xi+1}$ при любом $\xi \geq 0$. Значит,
        \[0<\alpha-\delta_{2\xi}<\delta_{2\xi+1}-\delta_{w\xi}=\frac{1}{Q_{2\xi}Q_{2\xi+1}},\ \ 0<\delta_{2\xi+1}-\alpha<\delta_{2\xi+1}-\delta_{w\xi}=\frac{1}{Q_{2\xi}Q_{2\xi+1}}\]
        Пусть $\tau = \frac{1+\sqrt{5}}{2}>1,6$. Докажем, что $Q_{\xi}\geq \tau^{\xi-1}$\\
        Индукция по $\xi$. База: $Q_1=1=\tau^{1-1}$. Пусть доказано для всех $\xi: q\leq \xi \leq m$.\\
        $\tau^2=\tau+1 \Rightarrow \tau^{k+2}=\tau^{k+1}+\tau^k$ для всех $k\geq 0$.\\
        Тогда
        \[Q_{m+1}=q_{m+1}Q_m+Q_{m-1}\geq Q_m+Q_{m-1}\geq \tau^{m-1}+\tau^{m-2}=\tau^m\]
        \begin{multline*}
            Q_{2\xi}Q_{2\xi-1}\geq \tau^{2\xi-1}\tau^{2\xi-2}=\tau^{4\xi-3}=\frac{\tau^{4\xi}}{\tau^3},\ 0<\alpha-\delta_{2\xi},\ \delta_{2\xi+1}-\alpha<\frac{\tau^{3}}{\tau^{4\xi}}=\\
            =\frac{2+\sqrt{5}}{(\frac{7+3\sqrt{5}}{2})^{\xi}}<\frac{5}{6^{\xi}}\to 0\ \ \ (\xi\to +\infty)
        \end{multline*}
    \end{proof} 
    \begin{comm}
        Последовательность неполных частных $q_0, q_1, q_2, \dots$ (бесконечная) периодична (начиная с некоторого номера) $\lra \alpha$ - квадратичная иррациональность, то есть $\alpha=\frac{A+B\sqrt{D}}{C},\ A,B,C,D\in \Z,\ D\geq 1$ - бесквадратное.
        Например:
        \[\sqrt{2}=[1;2,2,\dots]=1+\cfrac{1}{2+\cfrac{1}{2+\cfrac{1}{\ddots}}}\] 
    \end{comm} 
    \begin{comm}
        Лишь для немногих $\alpha\not\in \Q$, не являющихся квадратичной иррациональностью, известны разложения в цепную дробь.\\
        Число $e=2,718281828459045\dots$
        \[\alpha=\cfrac{e-1}{e+1}=[0;2,6,10,14,18,22,\dots]\]
        \[\alpha=\cfrac{e^{\frac{2}{k}}-1}{e^{\frac{2}{k}}+1}=\th{(e^{\frac{1}{k}})}=[0;1k,3k,5k,7k,\dots]\ \ \ \ \text{(Л. Эйлер, 1737г.)}\]
        \[e=[2;1,2,1,1,4,1,1,6,1,1,8,\dots]\ \ \ \ \text{(Л. Эйлер, 1737г.)}\]
        \[\alpha=\sqrt[3]{2}=[1;3,1,5,1,1,4,1,1,8,\dots]\]
        \[\alpha=\sqrt[3]{6}=[1;1,4,2,7,3,508,1,5,5,\dots]\]
        \[\pi=[3;7,5,1,292,1,1,1,2,1,\dots]\]
    \end{comm}
    \begin{comm}
        Рассмотрим $N\geq 3$ и дроби $\frac{a}{N},\ 1\leq a\leq N-1,\ (a,N)=1$ (таких дробей $\phi(N)$ штук). Разложим каждую в непрерывную дробь:
        \[\frac{a}{N}=\cfrac{1}{q_1+\cfrac{1}{q_2+\cfrac{1}{\ddots+\cfrac{1}{q_n}}}}\]\
        где $n=n(a)$ - длина разложения.
        Вопрос: Каково среднее значение $n(a)$ при изменении $a$?
        \[\frac{1}{\phi(N)}\sum\limits_{a=1}^N n(a)\approx \frac{12}{\pi^2}(\ln{2})+A,\ \ \text{A - некоторая константа}\ \ \ \text{(Х. Хейльбрин, 1969г.)}\]
    \end{comm}
    \section{Теория сравнений}
    Пусть $m\geq 2$ - целое.
    \begin{definition}
        Целые числа $a$ и $b$ называются сравнимыми по модулю $m$, если $a-b$ делится на $m$. Или, что то же, когда $a$ и $b$ при делении на $m$ дают одинаковые остатки. Число $m$ при этом называется модулем сравнения. Пишут $a\equiv b \pmod{m}$.
    \end{definition}
    \begin{examples}
        $8\equiv 2 \pmod{3},\ 15\equiv 1 \pmod{7},\ 24\equiv0 \pmod{6}$
    \end{examples}
    \begin{theorem} (Простейшие свойства сравнений)\label{th8.1}
         \begin{enumerate}
            \item $a\equiv a \pmod{m},\ \forall a$
            \item если $a\equiv b \pmod{m}$, то и $b\equiv a \pmod{m}$
            \item если $a\equiv b \pmod{m}, b\equiv c \pmod{m}$, то $a\equiv c \pmod{m}$
            \item если $a\equiv b \pmod{m},\ c\equiv d \pmod{m}$, то $a+c \equiv b+d \pmod{m}$ и\\
            $ac\equiv bd \pmod{m}$
         \end{enumerate}
    \end{theorem} 
    \begin{proof}
        1-3 очевидно. Докажем 4:\\
        $a=b+km,\ c=d+lm \Rightarrow a+c=(b+d)+(k+l)m\equiv b+d \pmod{m}$\\
        $ac=(b+km)(d+lm)=bd+blm+dkm+klm^2=bd+m(bl+dk+lkm)\equiv \\ \tab[13.5cm] \equiv bd \pmod{m}$
    \end{proof} 
    \begin{consequense}
        $a\equiv b \pmod{m},\ c,n$ - целые числа, $n\geq 1$, то $ca^n\equiv cb^n \pmod{m}$
    \end{consequense} 
    \begin{consequense}
        $a\equiv b \pmod{m},\ P(x)$ - многочлен с целыми коэффициэнтами\\ $\Rightarrow P(a)\equiv P(b) \pmod{m}$.
    \end{consequense} 
    \begin{theorem}\label{th8.2}
        Пусть $m\geq 2$. Тогда 
        \begin{enumerate}
            \item Если $ab\equiv ac \pmod{m}$ и $(a,m)=1$, то $b\equiv c \pmod{m}$.
            \item Если $a\equiv b \pmod{m},\ d\geq 1$ - целое, то $ad\equiv bd \pmod{md}$.
            \item Если $a\equiv b \pmod{m}$, то $(a,m)=(b,m)$.
            \item Если $a\equiv b \pmod{m}$ и $d \div m,\ d\geq 2$, то $a\equiv b \pmod{d}$.
        \end{enumerate}
    \end{theorem} 
    \begin{proof}\tab
        \begin{enumerate}
            \item $ab\equiv ac+km \Rightarrow a(b-c)=km$, если $b=c$, то утверждение очевидно. Пусть $b\ne c \Rightarrow b-c=0$ и из равенства $a(b-c)=km$ следует, что $a\div km$. Но $(a,m)=1 \Rightarrow$ по теореме \ref{th2.3} $a\div k$, то есть $k=an \Rightarrow a(b-c)=anm\\
            \Rightarrow b-c=nm \Rightarrow b\equiv c \pmod{m}$.
            \item $a=b+km \Rightarrow ad=bd+kmd \Rightarrow ad\equiv bd \pmod{md}$.
            \item По лемме \ref{lemma3.2} если $\alpha=bq+r$, то $(\alpha,\beta)=(\beta,r)$. Тогда $a\equiv b \pmod{m} \Rightarrow a=b+km \Rightarrow (\alpha=a,\ \beta=m,\ r=b) \Rightarrow (\alpha,\beta)=(a,m)=(\beta,r)=(m,b)$.
            \item $a=b+km,\ m=nd,\ d\geq 2 \Rightarrow a=b+knd \Rightarrow a\equiv b \pmod{d}$.
        \end{enumerate}
    \end{proof} 
    \begin{examples}\tab
        \begin{enumerate}
            \item $27 \equiv 3 \pmod{4} \Rightarrow 3\equiv 1 \pmod{4}$.
            \item $26\equiv 4 \pmod{22} \Rightarrow 13\equiv 2 \pmod{11}$.
            \item $48\equiv 28 \pmod{10}$, но $12 \not\equiv 7 \pmod{10}$.
        \end{enumerate}
    \end{examples}
    \begin{example}
        Найти остаток от деления $11^6$ на $9$.\\
        $11\equiv 2 \pmod{9} \Rightarrow  11^6 \equiv 2^6 \pmod{9}\Rightarrow 11^6\equiv (2^3)^2 \pmod{9}\\
        \Rightarrow 11\equiv (-1)^2 \pmod{9} \Rightarrow 11\equiv 1 \pmod{9}$.
    \end{example}
    \begin{example}
        Натуральное число $n=8k+7$ не представимо суммой трех квадратов целых чисел.\\
        $x$ - четное $\Rightarrow x=4y$ либо $x=4y+2,\ x^2=(4y)^2=16y^2\equiv 0 \pmod{8}$.\\
        $x^2=(4y+2)^2=16y^2+16y+4\equiv 4 \pmod{8}$.\\
        $x$ - нечетное $\Rightarrow x=4y\pm 1 \Rightarrow 16y^2\pm 8y+1\equiv 1 \pmod{8}$.
        Следовательно, $x^2\equiv 0,1,4 \pmod{8}$, но число $7$ нельзя представить суммой трех величин, принимающих значения $0$, $1$ и $4$. 
    \end{example}
    \begin{theorem*} (Теорема Лагранжа)\\
        Если $n\ne 4^a(8k+7)$, то $n$ представимо суммой трех квадратов целых чисел.
    \end{theorem*} 
    Сколько может быть чисел $n: 1\leq n\leq x,\ x\to +\infty,\ n=a^2+b^2$\ ?\\
    Таких чисел примерно $\cfrac{Bx}{\sqrt{\ln{x}}}$, где $B=0,7\dots$ - постоянная Рамануджана-Ландау.
    \begin{theorem}(8.4???)\label{th8.3}
        Если числа $a$ и $b$ сравнимы по модулям $m_1,\dots,m_k$, то они сравнимы по модулю $m=$НОК$(m_1,\dots,m_k)$. 
    \end{theorem} 
    \begin{proof}
        Если $a=b$, то утверждение очевидно. Пусть $a\ne b$, не теряя общности будем считать, что $a>b$. Так как $a\equiv b \pmod{m_i},\ i\in\overline{1,k} \Rightarrow$ натуральное число $a-b\equiv 0 \pmod{m}$, то есть число $a-b$ делится на каждое из чисел $m_1,\dots,m_k \Rightarrow a-b$ - их общее кратное. По теореме \ref{th2.1} получаем, что\\
        $[m_1,\dots,m_k] \div (a-b)$
    \end{proof} 
    Пусть задан модуль $m\geq 2$. Все множество $\Z$ разобьем на непересекающиеся подмножества, относя к одному и тому же подмножеству те числа, что при делении на $m$ дают дают одинаковые остатки. Именно, $a=q_1 m+r_1,\\
    b=q_2 m+r_2,\ 0\leq r_1,\ r_2\leq m-1$  относятся к одному и тому же подмножеству $\lra r_1=r_2$. Так получим ровно $m$ подмножеств, которые отвечают остаткам $r=0,1,\dots,m-1$ (все они непусты).
    \begin{definition}
        Построенные таким образом подмножества $\N$ называются классами вычетов по модулю $m$. Элементы каждого из этих подмножеств называются вычетами этого класса. Класс вычетов по модулю $m$, содержащий число $a$, иногда обозначают через $\bar{a}$ или $[a]$ или $[a]_m$.
    \end{definition} 
    Очевидно, что равенство классов $\bar{a}$ и $\bar{b}$ имеет место $\lra a\equiv b \pmod{m}$. Множество всех классов вычетов по модулю $m$ будем обозначать символом $\Z_m$
    \begin{example}
        Пусть $m=4 \Rightarrow$ остатки: $0,1,2,3 \Rightarrow \Z_4$ состоит из классов:\\
        $1)\ a=4n,\ \ \ \ 2)\ a=4n+1,\ \ \ \ 3)\ a=4n+2,\ \ \ \ 4)\ a=4n+3$.
    \end{example}
    \begin{definition}
        Пусть $m\geq 2$ и пусть $a_1,\dots,a_m$ - произвольные представители различных классов вычетов по модулю $m$. Тогда совокупность $a_1,\dots,a_m$ называется полной системой вычетов по модулю $m$.
    \end{definition} 
    \begin{examples}\tab
        \begin{enumerate}
            \item $m=4$, числа $a_1=13,\ a_2=7,\ a_3=6,\ a_4=8$.\\
            $a_4\equiv 0 \pmod{4},\ a_1\equiv 1 \pmod{4},\ a_2\equiv 2 \pmod{4},\ a_3\equiv 3 \pmod{4}$\\
            $\Rightarrow 13,7,6,8$ - полная система вычетов по модулю $4$.
            \item $m=5,\ a_1=2,\ a_2=6\ a_3=16,\ a_4=8,\ a_5=9$.\\
            $a_1 \equiv 2 \pmod{5},\ a_2 \equiv 1 \pmod{5},\ a_3 \equiv 1 \pmod{5},\ a_4 \equiv 3 \pmod{5},\\
            a_5 \equiv 4 \pmod{5} \Rightarrow$ числа $2,6,16,8,9$ не образуют полную систему вычетов по модулю $5$.
        \end{enumerate}
    \end{examples}
    Обычно в качестве полной системы вычетов по модулю $m$ берут совокупность $0,1,\dots, m-1$, состоящую из наименьших неотрицательных представителей всех классов вычетов.\\
    Иногда удобно работать с системой вычетов, составленой из наименьших по абсолютной величине представителей классов вычетов.
    \begin{example}
        Пусть $m=7$:\\
        $...,-21,\ -14,\ -7,\ 0,\ 7,\ 14,\ 21,... \equiv 0 \pmod{7}$ - берем 0\\
        $...,-20,\ -13,\ -6,\ 1,\ 8,\ 15,\ 22,... \equiv 1 \pmod{7}$ - берем 1\\
        $...,-19,\ -12,\ -5,\ 2,\ 9,\ 16,\ 23,... \equiv 2 \pmod{7}$ - берем 2\\
        $...,-18,\ -11,\ -4,\ 3,\ 10,\ 17,\ 24,... \equiv 3 \pmod{7}$ - берем 3\\
        $...,-17,\ -10,\ -3,\ 4,\ 11,\ 18,\ 25,... \equiv 4 \pmod{7}$ - берем -3\\
        $...,-16,\ -9,\ -2,\ 5,\ 12,\ 19,\ 26,... \equiv 5 \pmod{7}$ - берем -2\\
        $...,-15,\ -8,\ -1,\ 6,\ 13,\ 20,\ 27,... \equiv 6 \pmod{7}$ - берем -1\\
        Итак, полная наименьшая по абсолютной величине система вычетов\\ 
        по модулю $7: \{-3,-2,-1,1,2,3,0\}$
    \end{example}
    \begin{example}
        Общий случай:\\
        Для нечетного $n$ получаем $-\frac{m-1}{2},-\frac{m-3}{2},\dots,-1,0,1,\dots,\frac{m-3}{2},\frac{m-1}{2}$.\\
        Для четного $n$ получаем $-\frac{m}{2}+1,\dots,-1,0,1,\dots,\frac{m}{2}-1,\frac{m}{2}$.
    \end{example}
    \begin{theorem} \label{th8.4}
        Пусть $m\geq 2,\ a,b\in \Z$, причем $(a,m)=1$. Если величина $x$ пробегает полную систему вычетов по модулю $m$, то и величина $ax+b$ пробегает полную систему вычетов по модулю $m$.
    \end{theorem} 
    \begin{proof}
        Достаточно доказать, что если $x_1\not\equiv x_2 \pmod{m}$, то сравнение $ax_1\equiv ax_2 \pmod{m}$ невозможно. По теореме \ref{th8.3} п.1 на $a$ можно сократить: получим $x_1 \equiv x_2 \pmod{m}$ - противоречие.
    \end{proof} 
    \begin{consequense}
        Пусть $m\geq 2,\ a\in \Z,\ (a,m)=1$. Тогда существует единственный класс вычетов $c \pmod{m}$ такой, что $ac\equiv 1 \pmod{m}$.
    \end{consequense} 
    \begin{proof}
        Для $ax-1$ при некотором $x=c$ будет выполнено: $ac-1\equiv 0 \pmod{m},\ ac\equiv 1 \pmod{m}$. Пусть $ac_1 \equiv 1 \pmod{m}$ и $ac_2 \equiv 1 \pmod{m} \Rightarrow a(c_1-c_2)\equiv 0 \pmod{m} \Rightarrow m\div a(c_1-c_2) \Rightarrow m\div (c_1-c_2)$, то есть это возможно лишь при $c_1\equiv c_2 \pmod{m}$.
    \end{proof}
    \begin{comm}
        Такой вычет $c$ (класс вычетов $\bar{c}$) называют обратным к $a$ (соответственно обратным к классу $\bar{a}$). Обозначим его как $a^{\ast}$ (соответственно $\bar{a}^{\ast}$).
    \end{comm} 
    \begin{example}
        $m=5,\ a=3,\ b=4$\\
        $x\tab 3x=4\tab 3x+4\pmod{5}\\
        0\tab[1cm] 4\tab[2.5cm] 4\\
        1\tab[1cm] 7\tab[2.5cm] 2\\
        2\tab[1cm] 10\tab[2.3cm] 0\\
        3\tab[1cm] 13\tab[2.3cm] 3\\
        4\tab[1cm] 16\tab[2.3cm] 1$
    \end{example}
    \begin{comm}
        Условие $(a,m)=1$ опустить нельзя.
    \end{comm} 
    \begin{example}
        ПРИМЕР ПОЧЕМУ НЕЛЬЗЯ ОПУСКАТЬ
    \end{example}
    \begin{example}
        $m=7\\
        1\cdot 1\equiv 1\pmod{7}\Rightarrow 1^{ast}\equiv 1\pmod{7}\\
        2\cdot 4 \equiv 1\pmod{7}\Rightarrow 2^{ast}\equiv 4 \pmod{7}\\
        3\cdot 5 \equiv 1\pmod{7}\Rightarrow 3^{\ast}\equiv 5 \pmod{7}\\
        4\cdot 2 \equiv 1\pmod{7} \Rightarrow 4^{\ast}\equiv 2\pmod{7}\\
        5\cdot 3\equiv 1\pmod{7} \Rightarrow 5^{\ast}\equiv 3\pmod{7}\\
        6\cdot 6 \equiv 1\pmod{7} \Rightarrow 6^{\ast}\equiv 7\pmod{7}$
    \end{example}
    Согласно теореме \ref{th8.3}\ (пункт 3), числа, принадлежащие одному классу вычетов по модулю $m$, имеют с модулем один и тот же НОД. ($a\equiv b \pmod{m}\\
    \Rightarrow (a,m)=(a,b)$)\\
    Поэтому особый интерес придставляют классы, для которых этот НОД равен 1. Взяв от каждого такого класса по одному вычету, получим приведенную систему вычетов по модулю $m$. Возьмем в качестве такой полной системы вычетов числа $0,1,\dots, m-1$. Так как среди этих чисел количество взаимнопростых с модулем $m$ равно $\phi(m)$, то и любая приведенная система вычетовсодержит $\phi(m)$ элементов. Обозначение: $\Z_m^*$.
    \begin{example}
        $m=6;\ 0,1,2,3,4,5\ \Rightarrow 1,5$ - приведенная система вычетов.\\
        $m=7; 0,1,2,3,4,5,6\ \Rightarrow 1,2,3,4,5,6$ - приведенная система вычетов.\\
        $m=10; 0,1,2,3,4,5,6,7,8,9\ \Rightarrow 1,3,7,9$ - приведенная система вычетов.\\
        $m$ - простое $\Rightarrow \Z_m^*=\{1,2,\dots, p-1\}$.
    \end{example}
    \begin{theorem}\label{th8.5} (8.6??)
        Пусть $m\geq 2$, $a$ - целыое число, $(a,m)=1$, и пусть $x$ пробегает приведенную систему вычетов по модулю $m$. Тогда и величина $ax$ будет пробегать приведенную систему вычетов по модулю $m$.  
    \end{theorem} 
    \begin{proof}
        Что нужно проверить.
        \begin{enumerate}
            \item $ax_1\equiv ax_2$ невозможно, если $x_1\not\equiv x_2\pmod(m)$.
            \item $(ax,m)=1$ для всех $x\in \Z_m^*$.
        \end{enumerate}
        1. был проверен при доказательстве теоремы \ref{th8.4}.\\
        2. пусть $(ax,m)=\delta>1 \Rightarrow$ для некоторого $x: (x,m)=1\Rightarrow \delta \div ax$, причем $a\ne 0$ и $x\ne 0$ (следует из взаимной простоты с $m$) $\Rightarrow$ (по теореме \ref{th2.3}) $\delta \div a$. Но $\delta \div m$. Значит $\delta \div (a,m)\Rightarrow (a,m)\geq \delta>1$ противоречие.
    \end{proof} 
    \begin{theorem}\label{th8.6} (Теорема Эйлера)\\
        Пусть $m\geq 2$, $a$ - целое, $(a,m)=1\Rightarrow a^{\phi(m)}\equiv 1\pmod{m}$.
    \end{theorem} 
    \begin{proof}
        Пусть $1=r_1<r_2<\dots<r_c<\dots<r_{m-1}$, $c=\phi(m)$ - приведенная система вычетов. Пусть $ar_k\equiv \rho_k \pmod{m}$, где $0<\rho<m$. Из теоремы \ref{8.5} следует, что $\rho_1,\dots \rho_k$ образуют перестановку чисел $r_1,\dots r_k$. Перемножим сравнения почленно: $a^cr_1,\dots,r_c\equiv\rho_1\dots\rho_c \pmod{m}$. Но $r_1,\dots r_c=\rho_1\dots\rho_c=R$ и число $R$ взаимно просто с $m$ (следует из теоремы \ref{th2.3}). По теореме \ref{th8.3} (пункт 1), обе части сравнения $a^cR\equiv R\pmod{m}\Rightarrow a^{\phi(m)}\equiv 1\pmod{m}$.
    \end{proof} 
    \begin{consequense} (Малая теорема Ферма)\\
        Пусть $p$ - простое число. Тогда при любом целом $a$ выполняется сравнение: $a^p\equiv a\pmod{p}$.
    \end{consequense} 
    \begin{proof}
        Если $p\div a$, то очевидно. Если $(a,p)=1$, то $a^{\phi(p)}\equiv 1\pmod{p}\\
        \Rightarrow a^{p-1}\equiv 1\pmod{p}\Rightarrow a^p\equiv a\pmod{p}$.
    \end{proof}
    \newpage
    \section{Сравнения с одним неизвестным}
    Пусть $f(x)=a_nx^n+\dots+a_1x+a_0$ - многочлен с целыми коэффициэнтами. Будем изучать сравнения вида $(1)\ f(x)\equiv 0\pmod{m}$. Если $a_n\not\equiv 0\pmod{m}$, то число $m$ называется степенью сравнения. Решить сравнение (1) - значит найти все целые числа $x$, ему удовлетворяющие. По следствию 2 из \ref{th8.1},\ $a\equiv b\pmod{m}\Rightarrow f(a)\equiv f(b)\pmod{m}$. Значит, если (1) удовлетворяет некоторое число $x$, и $x\equiv a\pmod{m}$, то (1) удовлетворяют все числа сравнимые с $a$ по модулю $m$. По этой причине весь класс вычетов $a \pmod{m}$ удобно считать за одно решение.
    \subsection{Сравнения первой степени}
    Всякое сравнение первой степени можно переписать в виде: $ax\equiv b\pmod{m}$. Рассмотрим сперва случай, когда $(a,m)=1$. По теореме \ref{th8.4} получаем, что такое сравнение имеет единственное решение.
    \begin{enumerate}
        \item Способ 1. По теореме Эйлера получим $x_0\equiv ba^{\phi(m)-1}\pmod{m}$.\\
        Тогда $ax_0\equiv aba^{\phi(m)-1}\equiv ba^{\phi(m)}\pmod{m}\equiv b\cdot 1 \pmod{m}\equiv b\pmod{m}$.
        \item Способ 2. (Разложение в непрерывную дробь)
        $\alpha=\frac{m}{a}$ подходящие дроби: $\delta_0,\delta_1,\dots,\delta_{s-1}=\frac{P_{s-1}}{Q_{s-1}}, \delta_s=\frac{P_s}{Q_s}=\frac{m}{a}$. Известно по лемме \ref{lemma7.1} $P_s Q_{s-1}-P_{s-1}Q_s=(-1)^{s-1}\Rightarrow mQ_{s-1}-aP_{s-1}=(-1)^{s-1}$. $aP_{s-1}=(-1)^s+mQ_{s-1}\Rightarrow a(-1)^s P_{s-1}=1+mQ_{s-1}(-1)^s\Rightarrow a(-1)^sP_{s-1}b=b+mbQ_{s-1}(-1)^s$. Переходя к сравнению по модулю $m$, получим: $a(-1)^s P_{s-1}b=b \pmod{m}\Rightarrow (-1)^sP_{s-1}b\pmod{m}$ - решение.
    \end{enumerate}
    \begin{example}
        \[\frac{13}{7}=1+\frac{6}{7}=1+\cfrac{1}{\cfrac{7}{6}}=1+\cfrac{1}{1+\cfrac{1}{6}}\]
        Значит $x\equiv (-1)^2\cdot 2\cdot 3 \pmod{13}\equiv 6\pmod{13}$.
    \end{example}
    Пусть $(a,m)=d>1,\ ax\equiv b\pmod{m}$. Необходимое условие разрешимости - делимость $b$ на $d$.
    (т.к. если сравнение разрешимо, то $ax=b+km$ для некоторого целого $k$). Покажем что это условие достаточное. Пусть $a=a_1d,\ b=b_1d,\ m=m_1d,\ (a_1,m_1)=1$. Значит $a_1dx\equiv b_1d\pmod{m_1d}$. По теореме \ref{th8.3}, можно все сократить на $d: a_1x\equiv b_1\pmod{m_1}$. По доказаному выше, это сравнение имеет единственное решение по модулю $m_1$: $x\equiv x_1\pmod{m_1}$. Все числа вида (2)\ $x_1, x_1\pm m_1, x_1\pm 2m_1, \dots, x_1\pm tm_1,\dots$ - решения исходного сравненияю. Так как $x_1$ - решение, то $a_1x_1=b_1+km_1, k$ - некоторое целое число $\Rightarrow a(x_1\pm tm_1)=ax_1\pm tam_1=\alpha a_1 x_1 \pm t a_1 \alpha m_1 \equiv \alpha a_1 m_1 \pmod{m}\equiv \alpha(b_1+km_1)\pmod{m}\equiv b+km \pmod{m}\equiv b\pmod{m} \Rightarrow$ из ряда (2) нужно отобрать числа, различные по модулю $m$. $x_1,x_1+m_1,x_1+2m_1,\dots,x_1+(\alpha-1)m_1$ все они различны по модулю m.
    \begin{theorem}\label{th9.1}
        Пусть $m\geq 2, a,b$ - целые числа, причем $(a,m)=d$. Сравнение $ax\equiv b \pmod{m}$ разрешимо $\Leftrightarrow d\div b$. В случае разрешимости сравнение имеет $d$ решений.
    \end{theorem} 
    \begin{example}
        ПРИМЕР 9.2    
    \end{example}
\end{document}